\documentclass{article}

%===============================
%
%          📦 Paquetes
%
%===============================

% === Configuración del documento y tipografía ===
\usepackage[a4paper, top=2cm, bottom=2cm, left=2cm, right=2cm]{geometry}
\usepackage[spanish]{babel}
\usepackage[utf8]{inputenc}
\usepackage[T1]{fontenc}
\usepackage{eulervm}



% === Gráficos y representaciones visuales ===
\usepackage{tikz}
\usepackage{pgfplots}
\pgfplotsset{compat=1.18}
\usepackage{graphicx}
\usepackage{subcaption}
\usetikzlibrary{mindmap,trees}

% === Estilo y personalización ===
\usepackage{titling}
\usepackage{titlesec}
\usepackage{tocloft} 
\usepackage{fancyhdr}
\usepackage{setspace}
\usepackage{caption}

% === Matemáticas y símbolos ===
\usepackage{amsmath}
\usepackage{amssymb}
\usepackage{cancel}

% === Estructura y disposición del contenido ===
\usepackage{array}
\usepackage{multicol}
\usepackage{float}
\usepackage{indentfirst}

% === Hipervínculos y referencias ===
\usepackage{tocbibind}
\usepackage[colorlinks=true, linkcolor=black]{hyperref}

%===============================
%
%          ⬆️⬇️ Cabeceras
%
%===============================
\pagestyle{fancy}
\fancyhf{}
\fancyhead[L]{\includegraphics[width=3cm,clip,trim=1cm 1cm 1cm 1cm]{assets/logo-utp.png}}
\vspace{1mm}\fancyhead[R]{\textit{Huatay Salcedo, Luis Elías}\vspace{1mm}}
\fancyfoot[R]{\thepage}
\setlength{\headheight}{22.9188pt}
\addtolength{\topmargin}{-5.64944pt}

%======================================
%
%          📑 Espaciado de Párrafos
%
%======================================

\setlength{\parskip}{0.5em}


%======================================
%
%          📕 Estrucura del título
%
\title{
  \pagenumbering{gobble}
  \vspace{1cm}
  % \includegraphics[width=5cm,clip,trim=6cm 6cm 6cm 6cm]{./assets/isotipo.jpg} \\
  \includegraphics[width=5cm,clip,trim=0cm 2.9cm 0cm 0cm]{./assets/isotipo-utp.png} \\
  \vspace{0.5cm}
  \textbf{Universidad Tecnológica del Perú} \\
  \vspace{0.5cm}
  \text{Gestión de Proyectos} \\
  \vspace{1cm}
    {\huge \textbf{Tema 03: Tarea - Elaboración del EDT y Enunciado del Alcance del Proyecto}} \\
  \vspace{1cm}
  % \large \textbf{Para acceder al título de...} \\
}
\author{
  \begin{tabular}{ll}
    \textbf{Huatay Salcedo, Luis Elías, \texttt{U71809640}}\\
  \end{tabular} \\
  % \texttt{Ingeniería}\\
}
\begin{document}
% \date{14 de agosto de 2024}
\maketitle
\begin{center}
  Docente: \textbf{Rogelio Chambi Oscco}
\end{center}

%======================================
%
%          📚 Inicio del documento
%
%======================================

\newpage

% \begin{center}
% \Large\bfseries Índice
% \end{center}
% \begingroup
% \let\clearpage\relax
% \renewcommand{\contentsname}{}
% \tableofcontents
% \endgroup


\newpage
% Inicio de la numeración de páginas (desde aquí)
\pagenumbering{arabic}
\setcounter{page}{1}  
\vspace*{\fill}

\section{Introducción}

  El presente proyecto tiene como finalidad aplicar la Metodología de Gestión de Proyectos Clásica, siguiendo el marco de trabajo propuesto por el Project Management Institute (PMI) y utilizando herramientas del Framework del PMBOK. Durante la primera parte de la unidad, se desarrollará un marco predictivo para la gestión de un proyecto empresarial, abarcando las fases de Iniciación y Planificación.

  Para ello, se seleccionará una empresa y área específica, describiendo su misión, visión y objetivos estratégicos, así como el alineamiento del proyecto con dichos objetivos. Posteriormente, se documentarán los procesos de iniciación (objetivos, alcance inicial, gestión de interesados y acta de constitución) y planificación (gestión de integración, alcance, cronograma, costes, calidad, recursos, comunicaciones, riesgos y adquisiciones).

  Finalmente, se presentará un informe detallado que incluya todos los elementos mencionados, demostrando la capacidad de aplicar la metodología de gestión de proyectos en un entorno empresarial real.

\vspace*{\fill}

\section{Requerimientos del Proyecto}
\textit{Guía del PMBOK®}
\begin{itemize}
  \item \textbf{Requerimientos de negocio.}
  \begin{itemize}
    \item El proyecto debe contribuir a la transformación digital de los clientes a través de soluciones GIS.
    \item Debe permitir a las organizaciones optimizar procesos y mejorar la toma de decisiones mediante el análisis geoespacial.
    \item La consultoría debe generar valor estratégico en sectores empresarial, industrial y estatal.
    \item Posicionar a la empresa consultora como referente en el mercado de soluciones GIS.
  \end{itemize}

  \item \textbf{Requerimientos de los interesados.}
  \begin{itemize}
    \item Los clientes requieren reportes ejecutivos claros y comprensibles para la toma de decisiones.
    \item Los usuarios finales deben recibir capacitaciones básicas en el uso de la solución GIS.
    \item Los directivos esperan indicadores de desempeño e impacto que muestren beneficios tangibles.
    \item Los equipos de TI de los clientes necesitan que la solución sea compatible con la infraestructura existente.
  \end{itemize}

  \item \textbf{Requerimientos de la solución.}
  \begin{itemize}
    \item La solución debe integrar plataformas GIS licenciadas y/o de código abierto.
    \item Debe permitir la visualización, análisis y generación de mapas temáticos.
    \item La solución debe incluir módulos de carga, procesamiento y consulta de datos espaciales.
    \item Debe incorporar funcionalidades de interoperabilidad con sistemas existentes (ERP, BI, bases de datos).
    \item Documentación técnica completa y manuales de usuario.
  \end{itemize}

  \item \textbf{Requerimientos de transición.}
  \begin{itemize}
    \item Migración de datos espaciales desde fuentes externas o sistemas legados a la nueva plataforma GIS.
    \item Soporte técnico durante los primeros 3 meses posteriores a la implementación.
  \end{itemize}

  \item \textbf{Requerimientos del proyecto.}
  \begin{itemize}
    \item El proyecto debe ejecutarse en un plazo máximo de 8 semanas por cliente.
    \item El equipo de consultoría debe estar conformado por especialistas en GIS, programación y gestión de proyectos.
    \item Se debe elaborar un plan de gestión de riesgos y un plan de comunicaciones.
  \end{itemize}

  \item \textbf{Requerimientos de calidad.}
  \begin{itemize}
    \item La solución GIS debe garantizar una disponibilidad mínima del 95\% durante el primer año de uso.
    \item Los entregables deben cumplir con estándares internacionales de geoinformación (ej. ISO 19115, OGC).
    \item El nivel de satisfacción del cliente medido en encuestas de cierre debe superar el 90\%.
  \end{itemize}
\end{itemize}

\newpage

\subsection*{Necesidad}
Implementar una PMO

\subsection*{Requisitos de Negocio}
\begin{tabular}{|c|p{12cm}|}
\hline
\textbf{ID} & \textbf{Descripción} \\
\hline
RN1 & Gestionar los Proyectos de desarrollo de Software de modo Rentable \\
\hline
RN2 & Satisfacer las Necesidades de los clientes de Proyectos de desarrollo de software Exitosos y visibles \\
\hline
RN3 & Cumplimiento de Estandares y regulación respecto de proyectos de desarrollo de Software \\
\hline
\end{tabular}

\subsection*{Requisitos de Interesado}
\begin{tabular}{|c|p{7.5cm}|c|c|}
\hline
\textbf{ID} & \textbf{Descripción} & \textbf{Propietario} & \textbf{Relación} \\
\hline
RI1 & Contar con registro de Costos (facturación) por proyecto para asegurar rentabilidad & PMO & RN1,RN2 \\
\hline
RI2 & Visibilidad para cliente de sus proyectos y su portafolio por cliente y tipo de servicio en todo momento & Clientes & RN1,RN2 \\
\hline
RI3 & Registro de horas dedicadas de miembros del equipo (jira) & PMO & RN1,RN2 \\
\hline
RI4 & Poder auditar las aprobaciones de propuesta de desarrollo de parte del cliente & Auditoría & RN3 \\
\hline
RI5 & Registro de documentos que permanezcan si cambia personal de Tecnología & PMO & RN1 \\
\hline
RI6 & Registro de Caso de Negocio & PMO & RN1 \\
\hline
RI7 & Seguimiento de Proyectos Predictivos, Ágiles e Híbridos & PMO & RN1 \\
\hline
RI8 & Registro de encuesta de Satisfacción del Cliente por proyecto & PMO & RN2 \\
\hline
RI9 & Solución de software en Castellano & PMO & RN2 \\
\hline
\end{tabular}

\subsection*{Requisitos de Solución}

\subsubsection*{Requisitos Funcionales}
\begin{tabular}{|c|p{7.5cm}|c|c|}
\hline
\textbf{ID} & \textbf{Descripción} & \textbf{Propietario} & \textbf{Relación} \\
\hline
RSF1 & Registro de Costos (facturación) y visualización en todo momento de sus proyectos y su portafolio por cliente y tipo de servicio (Desarrollo, Infraestructura, Contabilidad) y del PM. & Clientes, PMO & RI1,RI2 \\
\hline
RSF2 & Cálculo automático de ROI e indicadores de avance & Clientes, PMO & RI2 \\
\hline
RSF3 & Registro, visualización y cálculo con horas de dedicación ingresadas por equipo & Clientes, PMO & RI3 \\
\hline
RSF4 & Visualización de proceso de aprobación de proyecto & Auditoría & RI4 \\
\hline
RSF5 & Definición y visualización de workflow de aprobación de proyectos & Clientes, PMO & RI5 \\
\hline
RSF6 & Registro y visualización de Caso de Negocio & Clientes, PMO & RI6 \\
\hline
RSF7 & Gestión y visualización de Proyectos Predictivos, Ágiles e Híbridos & Clientes, PMO & RI7 \\
\hline
RSF8 & Registro de encuesta de Satisfacción del Cliente por proyecto & Clientes, PMO & RI8 \\
\hline
RSF9 & Solución en Castellano & Clientes, PMO & RI9 \\
\hline
\end{tabular}

\subsubsection*{Requisitos No Funcionales}
\begin{tabular}{|c|p{9cm}|c|c|}
\hline
\textbf{ID} & \textbf{Descripción} & \textbf{Propietario} & \textbf{Relación} \\
\hline
RSNF1 & Seguridad para que acceso de clientes solo a su portafolio y PM solo a proyectos suyos & Clientes, PMO & RN3 \\
\hline
\end{tabular}

\subsection*{Requisitos de Transición}
\begin{tabular}{|c|p{9cm}|c|c|}
\hline
\textbf{ID} & \textbf{Descripción} & \textbf{Propietario} & \textbf{Relación} \\
\hline
RT1 & Mostrar Procesos y uso de formatos & PMO & RN1 \\
\hline
\end{tabular}

\section{Enunciado del Alcance del Proyecto}

El producto principal consiste en el diseño y desarrollo de una plataforma integral de servicios y entregables técnicos, que facilite a los clientes la optimización de sus procesos de toma de decisiones mediante el uso de información geoespacial. El desarrollo se realizará conforme a estándares internacionales de calidad y buenas prácticas en gestión de proyectos y GIS.

\subsection{Entregables}

\begin{itemize}
  \item Diagnóstico de necesidades y levantamiento de requerimientos GIS de los clientes.
  \item Base de datos geoespacial estructurada y documentada.
  \item Mapas temáticos digitales y en formato impreso (por ejemplo: uso de suelo, cobertura vegetal, infraestructura, riesgo).
  \item Informes técnicos de análisis espacial (por ejemplo: modelamiento de escenarios, identificación de patrones territoriales, análisis multitemporal).
  \item Plataforma digital para consulta y visualización de información geográfica.
  \item Capacitación inicial a los usuarios finales en el uso de herramientas GIS y en la interpretación de productos cartográficos.
\end{itemize}

\subsection{Entregables de Gestión (según guía PMBOK)}

\begin{itemize}
  \item Acta de Constitución del Proyecto.
  \item Gestión del alcance del proyecto.
  \item Estructura de Desglose del Trabajo (EDT).
  \item Plan de Gestión del Cronograma.
\end{itemize}

\subsection{Criterios de Aceptación}

Los entregables del proyecto serán aceptados si cumplen con los siguientes criterios:

\begin{itemize}
  \item Correspondencia con los requisitos establecidos por los interesados y documentados en la fase de planificación.
  \item Validación de la calidad de los datos espaciales (precisión, consistencia, completitud).
  \item Entrega de productos cartográficos y bases de datos en los formatos solicitados por el cliente.
  \item Aprobación formal por parte del patrocinador del proyecto y de los principales interesados.
\end{itemize}

\subsection{Exclusiones del Proyecto}

Quedan fuera del alcance del presente proyecto las siguientes actividades:

\begin{itemize}
  \item Mantenimiento y actualización continua de la base de datos geoespacial una vez entregada al cliente.
  \item Adquisición de licencias de software propietario (estas serán asumidas directamente por el cliente en caso de requerirse).
  \item Soporte técnico posterior al periodo de capacitación inicial.
  \item Integración de la plataforma con sistemas externos no contemplados en los requisitos iniciales.
\end{itemize}

Estas exclusiones permiten delimitar claramente las responsabilidades del proyecto, evitando desviaciones y gestionando de manera efectiva las expectativas de los interesados.

\section{Estructura de Desglose de Trabajo (EDT)}
\begin{multicols}{2}
  \subsection*{Nivel 1: Proyecto}
\begin{itemize}
  \item Implementación de Servicios de Consultoría y Asesoría en Soluciones GIS Empresariales e Institucionales
\end{itemize}

\subsection*{Nivel 2 y 3: Ciclo de Vida y Entregables}
\begin{enumerate}
  \item \textbf{Inicio}
  \begin{itemize}
    \item Acta de constitución del proyecto
    \item Registro de interesados
    \item Reunión de Kick-off
  \end{itemize}
  \item \textbf{Requerimientos}
  \begin{itemize}
    \item Levantamiento de necesidades de clientes (empresas, instituciones, gobierno, industria)
    \item Documentación de requisitos GIS (datos, plataformas, interoperabilidad, visualización)
    \item Validación y aprobación del documento de requisitos
  \end{itemize}
  \item \textbf{Análisis}
  \begin{itemize}
    \item Diagnóstico de datos y recursos geoespaciales existentes
    \item Identificación de brechas y necesidades de información
    \item Informe de análisis preliminar
  \end{itemize}
  \item \textbf{Diseño de la Solución}
  \begin{itemize}
    \item Diseño de la arquitectura de la base de datos GIS
    \item Diseño de mapas temáticos, dashboards y reportes empresariales
    \item Diseño de plataforma tecnológica (WebGIS, Desktop GIS, integración con sistemas empresariales)
  \end{itemize}
  \item \textbf{Desarrollo / Construcción}
  \begin{itemize}
    \item Creación y estructuración de la base de datos geoespacial
    \item Generación de mapas temáticos y reportes especializados
    \item Configuración de la plataforma tecnológica GIS (servidores, clientes, aplicaciones)
    \item Documentación técnica y operativa
  \end{itemize}
  \item \textbf{Pruebas y Validación}
  \begin{itemize}
    \item Validación de calidad de datos geoespaciales
    \item Validación de mapas temáticos y reportes con los interesados
    \item Pruebas de integración y usabilidad de la plataforma GIS
    \item Aceptación formal de entregables
  \end{itemize}
  \item \textbf{Implementación y Capacitación}
  \begin{itemize}
    \item Entrega oficial de la base de datos y mapas GIS
    \item Capacitación a usuarios (empresariales, técnicos, institucionales)
    \item Manuales técnicos y de usuario final
  \end{itemize}
  \item \textbf{Cierre}
  \begin{itemize}
    \item Informe final del proyecto
    \item Registro de lecciones aprendidas
    \item Acta de cierre
  \end{itemize}
\end{enumerate}

\subsection*{Nivel 2: Entregables de Gestión (PMBOK)}
\begin{enumerate}
  \item \textbf{Planificación}
  \begin{itemize}
    \item Plan de gestión del alcance
    \item Plan de gestión del cronograma
    \item Plan de gestión de costos
    \item Plan de gestión de riesgos
    \item Plan de calidad
  \end{itemize}
  \item \textbf{Seguimiento y Control}
  \begin{itemize}
    \item Informes de avance
    \item Monitoreo de cronograma y presupuesto
    \item Validación del alcance
    \item Control de cambios
  \end{itemize}
\end{enumerate}
\end{multicols}

\section*{Cronograma del Proyecto}

\textbf{Proyecto:} Implementación de PMO para proyectos de Asesoría en GIS

\begin{enumerate}
  \item \textbf{Definición de los requerimientos}
  \begin{enumerate}
    \item Análisis de necesidades de gestión de proyectos GIS
    \item Documento de especificaciones de requerimientos
    \item Acta de conformidad
  \end{enumerate}
  \item \textbf{Análisis y diseño de la solución}
  \begin{enumerate}
    \item Documento de modelado de procesos de gestión GIS
    \item Documento de análisis de herramientas y metodologías
    \item Documento de diseño de la PMO GIS
    \item Documento de políticas de calidad y seguridad
  \end{enumerate}
  \item \textbf{Construcción}
  \begin{enumerate}
    \item Desarrollo de soporte tecnológico
    \begin{enumerate}
      \item Configuración del software de gestión de proyectos (ej. ProjectLibre/MS Project)
      \item Desarrollo de plantillas y formatos estandarizados
    \end{enumerate}
    \item Pruebas unitarias de herramientas implementadas
    \item Acta de conformidad
  \end{enumerate}
  \item \textbf{Pruebas}
  \begin{enumerate}
    \item Elaboración de casos de prueba en proyectos piloto
    \item Realización de pruebas funcionales de la PMO
    \item Validación de satisfacción de interesados
    \item Acta de conformidad
  \end{enumerate}
  \item \textbf{Ajustes}
  \begin{enumerate}
    \item Realización de ajustes finales en procesos y formatos
    \item Acta de conformidad
  \end{enumerate}
  \item \textbf{Implementación}
  \begin{enumerate}
    \item Documento de pase a producción de la PMO
    \item Verificación de la operatividad de la PMO GIS
    \item Capacitación al equipo de proyectos GIS
    \item Manuales de procesos y herramientas de la PMO
    \item Acta de conformidad
  \end{enumerate}
  \item \textbf{Entregables de gestión de proyectos}
  \begin{enumerate}
    \item Acta de constitución del proyecto y registro de interesados
    \item Reunión de Kick-off
    \begin{enumerate}
      \item Reunión de coordinación para preparación de presentación e invitados \hfill 1h -- 01/10/2025
      \item Elaborar la presentación del proyecto \hfill 1d -- 02/10/2025
      \item Realizar la sesión de Kick-off \hfill 1h -- 03/10/2025
      \item Elaborar y enviar el acta del Kick-off \hfill 1h -- 03/10/2025
    \end{enumerate}
    \item Desarrollo y aprobación del plan de gestión del proyecto
    \item Entregables de ejecución
    \item Entregables de seguimiento y control
    \item Entregables de cierre
  \end{enumerate}
\end{enumerate}
\end{document}